% 第6章 对称性破缺的多天线纳米光子响应
\chapter{对称性破缺的多天线纳米光子响应}
\label{chap:result3}

\section{引言}
\label{sec:r3-intro}

% 超原子/超表面的设计原理
% 对称性在纳米光子学中的重要性
% 制造容差对超表面性能的影响

\section{超原子结构与制备}
\label{sec:r3-structure}

% 四个狭缝的zig-zag排列
% 名义上的镜像对称性
% 制备工艺

\begin{figure}[htbp]
    \centering
    % 占位符：用户实验数据图
    \fbox{\parbox{0.8\linewidth}{\centering\vspace{3cm} 图6-1：超原子结构\\（实验数据待插入）\vspace{3cm}}}
    \caption{\label{fig:ch6-meta} zig-zag排列的四狭缝超原子。(a) SEM图像；(b) 几何设计示意图，标注名义对称轴。}
\end{figure}

\section{实验测量结果}
\label{sec:r3-results}

% 振幅图和相位图
% 对称性破缺的直观观察
% 与数值模拟的偏差

\section{对称性破缺的机制分析}
\label{sec:r3-symmetry}

% $\sim$\SI{10}{nm}级别的加工不完美
% 相位差达$\pi/3$的定量分析
% 对电磁能量流的重新导向

\begin{figure}[htbp]
    \centering
    % 占位符：用户实验数据图
    \fbox{\parbox{0.8\linewidth}{\centering\vspace{3cm} 图6-2：对称性破缺测量\\（实验数据待插入）\vspace{3cm}}}
    \caption{\label{fig:ch6-symmetry} 超原子的对称性破缺响应。(a) 振幅图，干涉极大位置相对于对称预测的偏移；(b) 四个位置的时域信号，展示相干但不对称的场振荡。}
\end{figure}

\section{对超表面设计的启示}
\label{sec:r3-implication}

% 加工容差的分析
% 鲁棒性设计策略
% 主动补偿的可能性

\section{本章小结}
\label{sec:r3-summary}

本章研究了四狭缝超原子中结构对称性对电磁响应的影响。实验发现，即使名义上具有镜像对称的结构，$\sim$\SI{10}{nm}级别的加工偏差也足以导致电磁能量流的定向偏转，相位差可达$\pi/3$。这一结果揭示了超表面设计中加工精度与功能确定性之间的定量关系。
